\section{Introduction}
Today, 3D triangular surface meshes are used for a wide variety of purposes. Due to their ability to model objects efficiently, accurately, and in a customizable manner, they are commonly utilized in modeling, manufacturing, computer graphics, and geometry-based simulations. A mesh is often modeled from a real-world object, and hence it is not uncommon to see defects. Ju defines \textit{geometric correctness} as a state in which a mesh is watertight and free from defects \cite {ju2009geometric_errors}. Mesh defects are geometric errors, which include, but are not limited to, gaps, holes, non-manifold elements, and self-intersections. 

In this study, we specifically investigate intersections, a common defect where multiple components (i.e., facets or edges) within one mesh, or between two meshes, intersect each other. This may occur when surface meshes are created from image segmentations in which some of the compartments are locally thinner than a single voxel. Such intersecting surfaces do not allow the construction of a volume mesh, which is required for finite element modeling (FEM), a common mesh-based simulation method.

To address this issue, tools have been developed that can fix intersections. While some methods produce geometrically correct meshes, they may delete crucial mesh components, or introduce new problems such as low-quality elements. There is a need to find an optimal method depending on how much modification and complexity we allow, considering the intended application of the repaired mesh. Therefore, in this study, we consider not only geometric correctness of the repaired meshes, but also the quality of their triangles, and their structural integrity with respect to the input meshes. 

\subsection{Related Work}
When reviewing past work on intersection repairs, one can see two main approaches. The first involves removing intersecting components, whether they are facets or edges, and then patching the resulting holes. This is considered the simplest, most intuitive, and computationally lightweight method. Attene introduced such an algorithm to convert a mesh with defects into a single watertight triangular mesh~\cite{attene2010lightweight}. The algorithm iteratively eliminates intersecting triangles and their neighbors, then patches the resulting gaps with new edges and facets, until no self-intersections are detected. The main issue with this algorithm is that disconnected components may be removed during the degeneracy removal process. 

Another common approach is to detect intersections and resolve them by adding extra edges and vertices along crossing lines for triangulations. Its main advantage is that it does not involve major deletion and preserves the original mesh as much as possible; thus significant changes in volume and shape are not expected. Zhu developed an algorithm following this principle~\cite{Zhu2019robustalgo}. In his “direct repair algorithm,” Attene achieved the same goal by subdividing triangles along their intersecting elements~\cite{attene2014directrepair}. Although there are some variations of this approach, they all perform reasonably well in terms of retaining the original mesh structure~\cite{cherchi2020floating, charton2021topology}. However, this method is likely to create narrow and skewed triangles near areas where intersections once existed. If the sole objective is to remove intersections, this may not be a significant concern. However, in cases where mesh quality is crucial, such as in FEM, this can be problematic, as poorly shaped elements can negatively affect stability, accuracy, and convergence of the simulation~\cite{nemade2020meshquality}.

While the approaches discussed so far are considered mesh-based (or surface-based) repair methods, one can also voxelize a mesh (i.e., create a 3D cubic grid that resembles the shape of the surface mesh) and then create a completely new mesh based on this voxel grid~\cite{ju2009geometric_errors, sfikas2022FoR}, known as voxel-based repair. Several well-known algorithms, such as marching cubes \cite{Lorensen1987marchingcubes}, can be used to create the desired mesh~\cite{ju2009geometric_errors, sfikas2022FoR}. After this modification, no defects should remain \cite{sfikas2022FoR}. However, this method also presents challenges related to mesh quality, largely determined by the voxel size. If the voxel size is too large, the mesh loses a significant amount of its original information. If it is too small, the mesh complexity increases dramatically, resulting in a substantial computational burden for simulations.

Our goal in this study is to evaluate existing mesh repair approaches and assess their performance in terms of geometric correctness (i.e., amount of defects), mesh quality (i.e., triangle shape), and structural integrity (i.e., preservation of geometric features). 
We consider a wide range of quality metrics and evaluate their merits for this purpose. Furthermore, we explore alternative methods for resolving intersections, such as local remeshing and detaching sub-meshes, and examine their suitability for different scenarios. We provide all implemented algorithms, evaluation metrics, visualization methods, and data through our GitHub page.
